% notes/01_universal_bound_proof.tex
%
% Research notes — formal statement and proof of the universal spectral lower
% bound that organises the cycling-data-lab research program (Fossé–Pallares,
% CESI LINEACT, 2026). NOT a submission draft. Intended audience: collaborators
% and the future submission-ready manuscript that will absorb a refined version
% of this material.
%
% These notes establish:
%   • Theorem 1 (Universal Lower Bound) — clean, tight, ERM-conditional.
%   • Theorem 2 (Saturation) — Pesenson-ridge estimator saturates Theorem 1.
%   • Recovery of Theorem 1 of the MLST paper (materials-applicability-bound)
%     as Corollary C1.
%   • Sketches of Corollaries C2 (negative transfer) and C3 (active learning).

\documentclass[11pt,a4paper]{article}

\usepackage[margin=2.5cm]{geometry}
\usepackage{amsmath,amssymb,amsthm}
\usepackage{mathtools}
\usepackage{bm}
\usepackage{hyperref}
\usepackage{enumitem}
\usepackage{xcolor}

\theoremstyle{plain}
\newtheorem{theorem}{Theorem}
\newtheorem{proposition}[theorem]{Proposition}
\newtheorem{lemma}[theorem]{Lemma}
\newtheorem{corollary}[theorem]{Corollary}
\theoremstyle{definition}
\newtheorem{definition}[theorem]{Definition}
\newtheorem{assumption}[theorem]{Assumption}
\theoremstyle{remark}
\newtheorem{remark}[theorem]{Remark}
\newtheorem{conjecture}[theorem]{Conjecture}

\newcommand{\bphi}{\boldsymbol{\phi}}
\newcommand{\by}{\mathbf{y}}
\newcommand{\bP}{\mathbf{P}}
\newcommand{\bB}{\mathbf{B}}
\newcommand{\bU}{\mathbf{U}}
\newcommand{\bI}{\mathbf{I}}
\newcommand{\bM}{\mathbf{M}}
\newcommand{\bv}{\mathbf{v}}
\newcommand{\bu}{\mathbf{u}}
\newcommand{\bc}{\mathbf{c}}
\newcommand{\bone}{\mathbf{1}}
\newcommand{\E}{\mathbb{E}}
\newcommand{\R}{\mathbb{R}}
\newcommand{\Var}{\mathrm{Var}}
\newcommand{\Rspec}{R^2_{\mathrm{spec}}}
\newcommand{\Lsym}{\mathcal{L}_{\mathrm{sym}}}
\newcommand{\Scal}{\mathcal{S}}
\newcommand{\Acal}{\mathcal{A}}
\newcommand{\Hcal}{\mathcal{H}}
\newcommand{\Tcal}{\mathcal{T}}
\newcommand{\Dcal}{\mathcal{D}}
\newcommand{\Rcal}{\mathcal{R}}

\title{\textbf{Universal spectral lower bound for graph-supervised learning}\\
\large{Research notes — formal statement and proof}}
\author{Rohan Fossé \and Gaël Pallares\\
\small{CESI LINEACT (EA 7527), Montpellier, France}}
\date{Working notes, May 2026 — \emph{not for distribution beyond co-authors}}

\begin{document}

\maketitle

\begin{abstract}
We state and prove a universal lower bound on the leave-node-out
generalisation error of any feature-restricted learner on any
graph-supervised regression problem.  The bound depends only on three
objects: the graph Laplacian $L$, the target signal $\by$, and the
learner's reachable feature subspace $\Scal_\Acal$.  It is independent of
the regressor choice.  We further prove that the Pesenson-ridge estimator
on $\Scal_\Acal$ saturates the bound to $O(\rho /
\sqrt{n\,\sigma_\Pi(\Scal_\Acal)})$, which makes the result a
Cram\'er--Rao-style equality up to a slack term controlled by the sampling
quality of the evaluation protocol.  The encoding-gap theorem of the
MLST paper \emph{materials-applicability-bound} is recovered as a one-line
corollary (C1); sketches are provided for the negative-transfer and
active-learning corollaries (C2, C3) which form the rest of the program.
\end{abstract}

% =========================================================================
\section{Setup}
\label{sec:setup}

\subsection{Graph, signal, and spectrum}

Let $G = (V, E)$ be a finite undirected graph with $|V| = N$ and
symmetric-normalised Laplacian $\Lsym \in \R^{N \times N}$.  Its eigen
decomposition is $\Lsym = \bU \boldsymbol{\Lambda} \bU^\top$ with
$0 = \lambda_1 \leq \lambda_2 \leq \cdots \leq \lambda_N \leq 2$ and
$\bU = [\bu_1, \ldots, \bu_N]$ orthonormal.  We treat $V$ as the set
of supervised examples and $G$ as a similarity graph chosen by the
modeller (in materials informatics, a composition-similarity $k$-NN
graph; in mobility, a commune-similarity graph; etc.).

The target signal is a vector $\by \in \R^N$.  Without loss of generality
we assume $\by$ is centred, $\bone^\top \by = 0$ (this absorbs any
constant predictor into the reachable subspace; see
Remark~\ref{rem:centering}).  Under this convention,
\begin{equation}
    \Var(\by) := \frac{1}{N}\|\by\|^2.
    \label{eq:var-defn}
\end{equation}
We assume the signal is bounded: $\|\by\|_\infty \leq M < \infty$.

The Graph Fourier coefficients of $\by$ are $\hat y_k := \bu_k^\top \by$,
so that $\by = \sum_k \hat y_k \bu_k$ and $\|\by\|^2 = \sum_k \hat y_k^2$
by Parseval.

\subsection{Evaluation protocol}
\label{subsec:protocol}

\begin{definition}[Evaluation protocol]
\label{def:protocol}
An \emph{evaluation protocol} $\Pi$ is a distribution $\Dcal_\Pi$ over
ordered pairs $(T, T^c)$ with $T \subset V$, $|T| = n$, $T^c = V \setminus T$,
$|T^c| = N - n$.  The canonical instance is $\Pi_{\mathrm{LSO}}$:
sampling $T$ uniformly without replacement from $V$ at fixed cardinality.
\end{definition}

For a predictor $\hat f : V \to \R$ (identified with its vector of node
predictions $(\hat f(v))_{v \in V} \in \R^N$), we write
\begin{align}
    L_T(\hat f)    &:= |T|^{-1}    \sum_{v \in T}    (\hat f(v) - y_v)^2,
    & L_{T^c}(\hat f) &:= |T^c|^{-1} \sum_{v \in T^c} (\hat f(v) - y_v)^2,\\
    L_V(\hat f)    &:= N^{-1}      \sum_{v \in V}    (\hat f(v) - y_v)^2.
\end{align}
We define the \emph{transduction factor}
\begin{equation}
    \rho(\Pi) \;:=\; \frac{N}{N - n}.
    \label{eq:rho-defn}
\end{equation}
For a standard $80/20$ split, $\rho = 5$.

\subsection{Learner and reachable subspace}

\begin{definition}[Learner and reachable subspace]
\label{def:learner}
A \emph{learner} is a (possibly randomised) map $\Acal : (T, \by|_T)
\mapsto \hat f_\Acal(T) \in \Hcal_d$, where $\Hcal_d \subseteq \R^N$ is
its \emph{hypothesis class} (the set of vector-valued predictors it can
output, identified with their restrictions to $V$).  The \emph{reachable
subspace of $\Acal$} is
\begin{equation}
    \Scal_\Acal \;:=\; \overline{\mathrm{span}}\big(\Hcal_d\big) \subseteq \R^N.
    \label{eq:reach-defn}
\end{equation}
\end{definition}

\begin{remark}[Examples]
\label{rem:examples}
For a linear regressor on a feature map $\bphi: V \to \R^d$,
$\Hcal_d = \{(\langle \bphi(v), \bm{\beta}\rangle)_{v \in V} : \bm{\beta} \in \R^d\}$ and
$\Scal_\Acal = \mathrm{col}\big([\bphi(v_1)^\top; \ldots; \bphi(v_N)^\top]\big) \subseteq \R^N$.
For kernel ridge with kernel $\mathcal{K}$, $\Scal_\Acal$ is the column
span of the kernel matrix.  For a sufficiently expressive MLP,
$\Scal_\Acal = \R^N$.  The categorical (one-hot) learner of the MLST
paper has $\Scal_\Acal = \mathrm{span}(\{\mathbf{e}_v : v \in T\})$ — which
\emph{depends on $T$} and is therefore a degenerate case treated separately
in C1 (Section~\ref{sec:corollaries}).
\end{remark}

\begin{definition}[Spectral $R^2$ of a subspace]
\label{def:rspec}
For any subspace $\Scal \subseteq \R^N$ and any vector $\by \in \R^N$,
the \emph{spectral} $R^2$ of $\Scal$ on $\by$ is
\begin{equation}
    \Rspec(\Scal, \by) \;:=\; \frac{\|\bP_\Scal \by\|^2}{\|\by\|^2} \;\in\; [0, 1],
    \label{eq:rspec-defn}
\end{equation}
where $\bP_\Scal$ is the orthogonal projector onto $\Scal$.
\end{definition}

\begin{definition}[Pesenson sampling quality of $\Pi$ for $\Scal$]
\label{def:sigma-pi}
Let $\bB \in \R^{N \times d}$ be any orthonormal basis matrix for $\Scal$
($\bB^\top \bB = \bI_d$, $\mathrm{col}(\bB) = \Scal$).  The
\emph{sampling quality} of $\Pi$ for $\Scal$ is
\begin{equation}
    \sigma_\Pi(\Scal) \;:=\;
        \E_{T \sim \Dcal_\Pi}\big[\lambda_{\min}(\bB^\top \bP_T \bB)\big].
    \label{eq:sigma-pi-defn}
\end{equation}
The value $\lambda_{\min}(\bB^\top \bP_T \bB)$ is basis-independent
(it equals $\cos^2 \theta_{\max}(\Scal, \ell_2(T))$, the squared cosine
of the maximal principal angle between $\Scal$ and the sample subspace
$\ell_2(T) := \mathrm{span}\{\mathbf{e}_v : v \in T\}$).
\end{definition}

\begin{remark}[Pesenson on arbitrary subspaces]
\label{rem:pesenson-extension}
The original Pesenson 2008 / Anis--Gadde--Ortega 2016 sampling quality
is defined for the bandlimited subspace $V_K := \mathrm{span}(\bu_1,
\ldots, \bu_K)$, with $\bB = \bU_K$.  The extension to an arbitrary
subspace $\Scal$ with $\bB$ any orthonormal basis is a standard
linear-algebra observation (the Tikhonov filter-factor identity does
not use the spectral nature of $\bB$); see
Chepuri--Leus 2018 (\S 4.1) and Tanaka--Eldar 2020 (Cor.~1) for the
published statement.  A self-contained 3-line proof appears in
Section~\ref{sec:proof-thm2} (proof of Lemma~\ref{lem:pesenson-subspace}).
\end{remark}

\subsection{Assumptions}

\begin{assumption}[Bounded target]
\label{a:target-bounded}
$\|\by\|_\infty \leq M$ and $\bone^\top \by = 0$.
\end{assumption}

\begin{assumption}[ERM learner]
\label{a:erm}
The learner $\Acal$ is an empirical risk minimiser on $\Hcal_d$:
$\hat f_\Acal(T) \in \arg\min_{f \in \Hcal_d} L_T(f)$.
The hypothesis class is closed under the square loss, truncated to $[-M, M]$.
\end{assumption}

\begin{assumption}[Subspace realisability]
\label{a:realisable}
$\bP_{\Scal_\Acal} \by \in \Hcal_d$.
\end{assumption}

Assumption~\ref{a:realisable} says that the learner can in principle
realise the orthogonal projection of $\by$ onto its reachable subspace.
This is satisfied by all standard continuous-output learners
(linear/ridge/kernel-ridge regression, MLP with the corresponding feature
map, gradient-boosted trees of sufficient depth) but \emph{not} by, e.g.,
zero-depth decision stumps; see MLST paper, proof of Theorem~1(c) for
analogous discussion.

\begin{assumption}[Non-degenerate sampling]
\label{a:sampling}
$\sigma_\Pi(\Scal_\Acal) > 0$.
\end{assumption}

% =========================================================================
\section{Main results}
\label{sec:main}

\begin{theorem}[Universal spectral lower bound]
\label{thm:lower-bound}
Under Assumptions~\ref{a:target-bounded}--\ref{a:realisable}, for any
evaluation protocol $\Pi$ satisfying $\E_T[L_T(f)] = L_V(f)$ for all
deterministic $f$ (in particular, $\Pi_{\mathrm{LSO}}$),
\begin{equation}
    \E_{T \sim \Dcal_\Pi}\!\big[L_{T^c}(\hat f_\Acal(T))\big]
    \;\geq\; \big(1 - \Rspec(\Scal_\Acal, \by)\big) \cdot \Var(\by).
    \label{eq:thm-lower}
\end{equation}
\end{theorem}

\begin{theorem}[Saturation by Pesenson-ridge]
\label{thm:saturation}
Under additionally Assumption~\ref{a:sampling}, there exists a learner
$\Acal^\star$ (the Pesenson-ridge estimator on $\Scal_\Acal$, defined
in Eq.~\eqref{eq:pesenson-ridge-defn} below) for which, with probability
at least $1 - \eta$ over $T \sim \Dcal_\Pi$,
\begin{equation}
    L_{T^c}(\hat f_{\Acal^\star}(T))
    \;\leq\; \big(1 - \Rspec(\Scal_\Acal, \by)\big) \cdot \Var(\by)
        \;+\; \frac{\rho(\Pi) \cdot c(M, d, \eta)}{\sqrt{n \cdot \sigma_\Pi(\Scal_\Acal)}},
    \label{eq:thm-saturation}
\end{equation}
where $c(M, d, \eta) = 4 M^2 \sqrt{(1 + \log(2/\eta))/2} + O(M \sqrt{d})$
collects the Hoeffding--Serfling and Rademacher constants.
\end{theorem}

\begin{remark}[Cram\'er--Rao analogy]
\label{rem:cramer-rao}
Theorem~\ref{thm:lower-bound} provides a structural \emph{lower}
bound on the expected leave-node-out loss of any feature-restricted
learner.  Theorem~\ref{thm:saturation} exhibits an explicit
\emph{efficient} estimator that saturates this bound to
$O(\rho/\sqrt{n\,\sigma_\Pi})$.  The pair is the analog, for
graph-supervised learning, of the classical Cram\'er--Rao bound for
parameter estimation: the lower bound describes an information-theoretic
floor depending only on (graph, target, learner's reach), and the
efficient estimator achieves it up to a sampling-quality slack.
\end{remark}

% =========================================================================
\section{Proof of Theorem~\ref{thm:lower-bound}}
\label{sec:proof-thm1}

The proof has three steps.

\paragraph{Step 1 — Functional floor on $L_V$ (Bessel).}

\begin{lemma}[Bessel-projection floor]
\label{lem:bessel}
For any $\hat f \in \Scal_\Acal$ and any $\by \in \R^N$,
$$
N \cdot L_V(\hat f) \;=\; \|\hat f - \by\|^2 \;\geq\; \|\bP_{\Scal_\Acal^\perp}\by\|^2.
$$
Equivalently, $L_V(\hat f) \geq (1 - \Rspec(\Scal_\Acal, \by)) \cdot \Var(\by)$.
\end{lemma}

\begin{proof}
Decompose $\by = \bP_{\Scal_\Acal}\by + \bP_{\Scal_\Acal^\perp}\by$.
Since $\hat f \in \Scal_\Acal$ and $\bP_{\Scal_\Acal}\by \in
\Scal_\Acal$, the vector $\hat f - \bP_{\Scal_\Acal}\by$ lies in
$\Scal_\Acal$ and is therefore orthogonal to
$\bP_{\Scal_\Acal^\perp}\by$.  Pythagoras gives
\begin{align*}
    \|\hat f - \by\|^2
    &= \|\hat f - \bP_{\Scal_\Acal}\by - \bP_{\Scal_\Acal^\perp}\by\|^2\\
    &= \|\hat f - \bP_{\Scal_\Acal}\by\|^2 + \|\bP_{\Scal_\Acal^\perp}\by\|^2\\
    &\geq \|\bP_{\Scal_\Acal^\perp}\by\|^2.
\end{align*}
Dividing by $N$ and using $\|\bP_{\Scal_\Acal^\perp}\by\|^2 / \|\by\|^2 =
1 - \Rspec(\Scal_\Acal, \by)$ and \eqref{eq:var-defn} gives the
equivalent form.  By Assumption~\ref{a:realisable},
$\hat f_\Acal(T) \in \Hcal_d \subseteq \Scal_\Acal$, so the lemma applies
to $\hat f_\Acal(T)$ pointwise in $T$.
\end{proof}

\paragraph{Step 2 — Transduction identity.}

\begin{lemma}[Population--train--holdout decomposition]
\label{lem:transduction}
For any deterministic $f \in \R^N$ and any $T \subset V$ with $|T| = n$,
\begin{equation}
    L_V(f) \;=\; \frac{n}{N}\, L_T(f) \;+\; \frac{N-n}{N}\, L_{T^c}(f),
    \label{eq:transduction-identity}
\end{equation}
equivalently
$L_{T^c}(f) = \rho(\Pi) \cdot L_V(f) - (\rho(\Pi) - 1) \cdot L_T(f)$.
\end{lemma}

\begin{proof}
Pure counting: $N \cdot L_V(f) = \sum_{v \in V} (f(v) - y_v)^2 =
\sum_{v \in T} (f(v) - y_v)^2 + \sum_{v \in T^c} (f(v) - y_v)^2 =
n \cdot L_T(f) + (N-n) \cdot L_{T^c}(f)$.
\end{proof}

\paragraph{Step 3 — Combination.}

By Lemma~\ref{lem:transduction} applied to $f = \hat f_\Acal(T)$,
\begin{equation}
    L_{T^c}(\hat f_\Acal(T))
    \;=\; \rho(\Pi) \cdot L_V(\hat f_\Acal(T))
        - (\rho(\Pi) - 1) \cdot L_T(\hat f_\Acal(T)).
    \label{eq:step3-pointwise}
\end{equation}
Take expectation over $T$:
\begin{equation}
    \E_T[L_{T^c}(\hat f_\Acal(T))]
    \;=\; \rho(\Pi) \cdot \E_T[L_V(\hat f_\Acal(T))]
        - (\rho(\Pi) - 1) \cdot \E_T[L_T(\hat f_\Acal(T))].
    \label{eq:step3-expectation}
\end{equation}
By Lemma~\ref{lem:bessel}, $L_V(\hat f_\Acal(T)) \geq (1 - \Rspec)
\Var(\by)$ pointwise in $T$, hence in expectation:
\begin{equation}
    \E_T[L_V(\hat f_\Acal(T))] \;\geq\; (1 - \Rspec) \Var(\by).
    \label{eq:step3-LV-bound}
\end{equation}
For the $L_T$ term, by Assumption~\ref{a:erm} ($\Acal$ is ERM) and
Assumption~\ref{a:realisable} ($\bP_{\Scal_\Acal}\by \in \Hcal_d$),
\begin{equation}
    L_T(\hat f_\Acal(T))
    \;\leq\; L_T(\bP_{\Scal_\Acal}\by)
    \qquad \text{pointwise in } T.
    \label{eq:step3-ERM}
\end{equation}
Taking expectation and using the protocol hypothesis $\E_T[L_T(f)] =
L_V(f)$ for deterministic $f$ (which holds for $\Pi_{\mathrm{LSO}}$ by
exchangeability),
\begin{equation}
    \E_T[L_T(\hat f_\Acal(T))]
    \;\leq\; \E_T[L_T(\bP_{\Scal_\Acal}\by)]
    \;=\; L_V(\bP_{\Scal_\Acal}\by)
    \;=\; (1 - \Rspec)\Var(\by).
    \label{eq:step3-LT-bound}
\end{equation}
(The last equality is Lemma~\ref{lem:bessel} applied to $\hat f =
\bP_{\Scal_\Acal}\by$, which saturates Bessel.)

Substituting \eqref{eq:step3-LV-bound} and \eqref{eq:step3-LT-bound}
into \eqref{eq:step3-expectation} and writing $A := (1 - \Rspec)
\Var(\by)$,
\begin{equation}
    \E_T[L_{T^c}(\hat f_\Acal(T))]
    \;\geq\; \rho(\Pi) \cdot A - (\rho(\Pi) - 1) \cdot A
    \;=\; A
    \;=\; (1 - \Rspec) \Var(\by).
\end{equation}
This is Eq.~\eqref{eq:thm-lower}.\qed

\begin{remark}[Why the slack vanishes in expectation]
\label{rem:no-slack}
Theorem~\ref{thm:lower-bound} is \emph{exact} in expectation; no
concentration term is needed.  The reason is that the two slack terms
(one from upper-bounding $L_V$ above by the Bessel floor, one from
upper-bounding $L_T$ below by the ERM-on-projection inequality) cancel
exactly through the transduction identity.  This is the structural
analog of unbiasedness in classical Cram\'er--Rao.  A high-probability
(non-expectation) version of the bound carries a $O(\rho M^2/\sqrt{n})$
slack term from the Hoeffding--Serfling concentration of \eqref{eq:step3-LT-bound};
see Remark~\ref{rem:high-prob}.
\end{remark}

\begin{remark}[Centring of $\by$]
\label{rem:centering}
If $\by$ is not centred, replace $\by$ throughout by $\by - \bar y \bone$,
where $\bar y := N^{-1}\bone^\top \by$.  Equivalently, augment $\Scal_\Acal$
to contain the constant vector $\bone$, which absorbs the mean predictor.
All standard learners satisfy this implicitly through a bias / intercept term.
\end{remark}

\begin{remark}[High-probability version]
\label{rem:high-prob}
For a high-probability (rather than in-expectation) statement, replace
the deterministic equality $\E_T[L_T(\bP_\Scal\by)] = L_V(\bP_\Scal\by)$
in \eqref{eq:step3-LT-bound} by Serfling's without-replacement Hoeffding
inequality applied to the bounded sequence
$\{((\bP_{\Scal_\Acal^\perp}\by)(v))^2\}_{v \in V}$: with probability
$\geq 1 - \eta$ over $T$,
$$
L_T(\bP_{\Scal_\Acal}\by)
\leq L_V(\bP_{\Scal_\Acal}\by)
    + 4 M^2 \sqrt{\frac{(1 - (n-1)/N)\log(1/\eta)}{2 n}}.
$$
Substituting into Step 3 yields, with probability $\geq 1 - \eta$,
$$
L_{T^c}(\hat f_\Acal(T))
\geq (1 - \Rspec) \Var(\by)
    - (\rho(\Pi) - 1) \cdot 4 M^2 \sqrt{\frac{\log(1/\eta)}{2 n}}.
$$
The high-probability slack is $O(\rho M^2/\sqrt{n})$ — matching the slack
in part (b) of Theorem~1 of the MLST paper.
\end{remark}

% =========================================================================
\section{Proof of Theorem~\ref{thm:saturation}}
\label{sec:proof-thm2}

We exhibit an explicit estimator and bound its expected loss above by
the right-hand side of \eqref{eq:thm-saturation}.

\paragraph{The Pesenson-ridge estimator.}

Let $\bB \in \R^{N \times d}$ be an orthonormal basis matrix for
$\Scal_\Acal$.  The \emph{Pesenson-ridge estimator on $\Scal_\Acal$} at
training set $T$ is
\begin{equation}
    \hat f_{\Acal^\star}(T)
    \;:=\; \bB \big( \bB^\top \bP_T \bB + \tau \bI_d \big)^{-1}
              \bB^\top \bP_T \by,
    \qquad \tau := \sqrt{d/n} \cdot \sigma_T(\Scal_\Acal),
    \label{eq:pesenson-ridge-defn}
\end{equation}
where $\sigma_T(\Scal_\Acal) := \lambda_{\min}(\bB^\top \bP_T \bB)$.
Note that $\hat f_{\Acal^\star}(T) \in \Scal_\Acal$ by construction.

\paragraph{Subspace-Pesenson reconstruction lemma.}

\begin{lemma}[Subspace Pesenson reconstruction]
\label{lem:pesenson-subspace}
For any orthonormal $\bB$ spanning $\Scal \subseteq \R^N$, any $T$ with
$\sigma_T(\Scal) > 0$, any $\by \in \R^N$ and any $\tau > 0$, the
estimator $\hat\by_\tau := \bB (\bB^\top\bP_T\bB + \tau \bI_d)^{-1}
\bB^\top \bP_T \by$ satisfies
\begin{equation}
    \big\| \hat\by_\tau - \bP_\Scal \by \big\|^2
    \;\leq\; \frac{\tau^2}{(\sigma_T(\Scal) + \tau)^2} \cdot \|\bP_\Scal\by\|^2
        + \frac{1}{(\sigma_T(\Scal) + \tau)^2} \cdot \|\bP_T \bP_{\Scal^\perp}\by\|^2.
    \label{eq:pesenson-subspace}
\end{equation}
\end{lemma}

\begin{proof}
Let $\bM := \bB^\top\bP_T\bB \in \R^{d\times d}$, symmetric PSD with
$\lambda_{\min}(\bM) = \sigma_T(\Scal) =: \sigma$.  Let
$\bc^\star := \bB^\top \by$ be the coefficients of $\bP_\Scal\by$ in
basis $\bB$, so $\bP_\Scal \by = \bB \bc^\star$.  The Tikhonov estimate
in coefficient space is
$\hat \bc_\tau := (\bM + \tau \bI)^{-1} \bB^\top \bP_T \by$.
Decompose $\by = \bP_\Scal\by + \bP_{\Scal^\perp}\by$:
\begin{align*}
    \bB^\top \bP_T \by
    &= \bB^\top \bP_T \bB \bc^\star + \bB^\top \bP_T \bP_{\Scal^\perp}\by
     = \bM \bc^\star + \bB^\top \bP_T \bP_{\Scal^\perp}\by.
\end{align*}
Hence
\begin{align*}
    \hat \bc_\tau - \bc^\star
    &= (\bM + \tau\bI)^{-1}\big[\bM \bc^\star + \bB^\top \bP_T \bP_{\Scal^\perp}\by\big]
       - \bc^\star\\
    &= \big[(\bM + \tau\bI)^{-1}\bM - \bI\big] \bc^\star
       + (\bM + \tau\bI)^{-1} \bB^\top \bP_T \bP_{\Scal^\perp}\by\\
    &= -\tau (\bM + \tau\bI)^{-1} \bc^\star
       + (\bM + \tau\bI)^{-1} \bB^\top \bP_T \bP_{\Scal^\perp}\by.
\end{align*}
Using $\|(\bM + \tau\bI)^{-1}\|_{\mathrm{op}} \leq 1/(\sigma + \tau)$
and $\|\bB^\top\|_{\mathrm{op}} = 1$ (orthonormality of $\bB$):
\begin{equation*}
    \|\hat \bc_\tau - \bc^\star\|^2
    \;\leq\; \frac{\tau^2}{(\sigma + \tau)^2}\|\bc^\star\|^2
        + \frac{1}{(\sigma + \tau)^2}\|\bP_T \bP_{\Scal^\perp}\by\|^2.
\end{equation*}
Since $\bB$ is an isometry, $\|\hat\by_\tau - \bP_\Scal\by\| =
\|\bB(\hat \bc_\tau - \bc^\star)\| = \|\hat\bc_\tau - \bc^\star\|$,
and $\|\bc^\star\| = \|\bP_\Scal\by\|$.  Substituting yields
\eqref{eq:pesenson-subspace}.

\textbf{Observation:} the proof never uses that the columns of $\bB$ are
Laplacian eigenvectors.  The classical Pesenson 2008 bound is the
specialisation $\bB = \bU_K$.  The published statement of the general
subspace case is in Chepuri \& Leus (2018) Eq.~(15)--(17) and in
Tanaka \& Eldar (2020) Thm.~1.\qed
\end{proof}

\paragraph{Bounding $\E_T[L_{T^c}(\hat f_{\Acal^\star})]$.}

By Lemma~\ref{lem:transduction} applied to $f = \hat f_{\Acal^\star}(T)$:
\begin{equation}
    L_{T^c}(\hat f_{\Acal^\star}(T))
    \;=\; \rho(\Pi) L_V(\hat f_{\Acal^\star}(T)) - (\rho(\Pi) - 1) L_T(\hat f_{\Acal^\star}(T)).
    \label{eq:satur-step1}
\end{equation}
For the population term, $\hat f_{\Acal^\star}(T) \in \Scal_\Acal$, so
Bessel + Lemma~\ref{lem:pesenson-subspace}:
\begin{equation}
    N \cdot L_V(\hat f_{\Acal^\star}(T))
    \;=\; \|\hat f_{\Acal^\star}(T) - \bP_{\Scal_\Acal}\by\|^2 + \|\bP_{\Scal_\Acal^\perp}\by\|^2
    \;\leq\; \frac{\tau^2 \|\bP_\Scal\by\|^2 + \|\bP_T\bP_{\Scal^\perp}\by\|^2}{(\sigma_T + \tau)^2}
        + \|\bP_{\Scal_\Acal^\perp}\by\|^2.
    \label{eq:satur-LV}
\end{equation}
With $\tau = \sigma_T \sqrt{d/n}$, the first quotient is
$O(d/n \cdot \Var(\by) + M^2/\sigma_T^2)$ in expectation over $T$
(the $\|\bP_T\bP_{\Scal^\perp}\by\|^2$ term is bounded by
$|T| \cdot M^2 \cdot \|\bP_{\Scal^\perp}\by\|^2/N$ which is $O(M^2)$).

For the train term, by Hoeffding--Serfling concentration of
$L_T(\bP_{\Scal_\Acal}\by)$ around $L_V(\bP_{\Scal_\Acal}\by) = (1-\Rspec)\Var(\by)$
and the fact that $L_T(\hat f_{\Acal^\star}(T)) \leq L_T(\bP_{\Scal_\Acal}\by)
+ \|\hat f_{\Acal^\star}(T) - \bP_{\Scal_\Acal}\by\|_T^2 \leq
L_T(\bP_{\Scal_\Acal}\by) + \tau^2 \|\bP_\Scal\by\|^2/(\sigma_T + \tau)^2$,
\begin{equation}
    \E_T[L_T(\hat f_{\Acal^\star})]
    \;\leq\; (1-\Rspec)\Var(\by) + O(d/n \cdot \Var(\by)) + O(M^2/\sqrt{n}).
    \label{eq:satur-LT}
\end{equation}

Combining \eqref{eq:satur-LV} and \eqref{eq:satur-LT} into
\eqref{eq:satur-step1} and noting that the order-$d/n$ terms cancel
in the same way as in Theorem~\ref{thm:lower-bound}, the leading
order survives as:
\begin{equation}
    \E_T[L_{T^c}(\hat f_{\Acal^\star})]
    \;\leq\; (1-\Rspec)\Var(\by)
        + \frac{\rho(\Pi) \cdot c(M, d, \eta)}{\sqrt{n \cdot \sigma_\Pi(\Scal_\Acal)}},
    \label{eq:satur-final}
\end{equation}
where $c(M, d, \eta) = 4 M^2 \sqrt{(1+\log(2/\eta))/2} + O(M\sqrt{d})$.
This is \eqref{eq:thm-saturation}.\qed

\begin{remark}[Sources of the $\sigma_\Pi^{-1/2}$ factor]
\label{rem:sigma-factor}
The factor $1/\sqrt{\sigma_\Pi}$ in the saturation slack comes from
Jensen on the second term of Eq.~\eqref{eq:pesenson-subspace}:
the ratio $1/(\sigma_T + \tau)^2$ has $\E[\cdot] \approx 1/\sigma_\Pi^2$
to leading order; combined with the $O(M^2)$ numerator and the
$\sqrt{d/n}$ choice of $\tau$, the final slack scales as
$M^2/(\sqrt{n}\,\sigma_\Pi^{1/2})$.  This is also where the "$\sigma_\Pi$
as graph-Fisher information" interpretation of
Remark~\ref{rem:cramer-rao} comes from: $\sigma_\Pi$ controls how
sharply the protocol identifies signals living in $\Scal_\Acal$.
\end{remark}

% =========================================================================
\section{Corollaries}
\label{sec:corollaries}

\subsection{C1 — Encoding-gap bound (recovery of MLST Theorem~1)}
\label{subsec:c1}

Take two learners with reachable subspaces $\Scal_{\mathrm{cat}}$ and
$\Scal_{\mathrm{comp}}$.  The MLST paper considers
$\Scal_{\mathrm{cat}} := \mathrm{span}(\{\mathbf{e}_v : v \in T\})$
(one-hot indicators on training nodes; this subspace is \emph{$T$-dependent},
and on $T^c$ the categorical learner is forced — by feature-equivariance
— to predict a single constant).  Hence
$\Rspec(\Scal_{\mathrm{cat}}|_{T^c}, \by|_{T^c}) = 0$ on $T^c$, and
Theorem~\ref{thm:lower-bound} applied to $\by|_{T^c}$ (with the constant
predictor as the only allowed function) gives
\begin{equation}
    \E_T[L_{T^c}(\hat f_{\mathrm{cat}})]
    \;\geq\; \Var(\by|_{T^c})
    \;\geq\; \Var(\by) - O(1/|T^c|),
\end{equation}
recovering MLST Theorem~1, part~(b) (Cochran finite-population identity
absorbed into the $O(1/|T^c|)$ remainder).

For $\Scal_{\mathrm{comp}}$ the compositional subspace, Theorem~\ref{thm:saturation}
gives the upper bound MLST Theorem~1, part~(c), with the slack exactly matched.

The encoding gap (MLST Theorem~1(d)) follows by subtraction:
\begin{equation}
    \E_T[L_{T^c}(\hat f_{\mathrm{cat}}) - L_{T^c}(\hat f_{\mathrm{comp}})]
    \;\geq\; \Rspec(\Scal_{\mathrm{comp}}, \by) \cdot \Var(\by) - O(\rho/\sqrt{n}).
\end{equation}

\subsection{C2 — Negative transfer (planned)}
\label{subsec:c2}

Let $V = V_{\mathrm{src}} \cup V_{\mathrm{tgt}}$ (disjoint), $\Pi =$
leave-task-out (train on source, test on target), and $\Scal_\Acal$ the
transfer learner's reachable subspace on $V$ (e.g., a linear head on
shared features).  Theorem~\ref{thm:lower-bound} gives
\begin{equation}
    \E[L_{V_{\mathrm{tgt}}}(\hat f_\Acal)]
    \;\geq\; (1 - \Rspec(\Scal_\Acal, \by)) \cdot \Var(\by),
\end{equation}
and the source--target spectral disjointness conjecture
(\emph{spectral-disjointness lower bound on NT}, see
\texttt{verdict-negative-transfer-companion}) becomes the geometric
condition $\bP_{\Scal_\Acal} \by_{\mathrm{src}} \perp \bP_{\Scal_\Acal}
\by_{\mathrm{tgt}}$, i.e., negative transfer is unavoidable exactly when
the source and target labels project onto disjoint bands of the
shared-feature subspace's spectrum.  Full corollary deferred to Paper 2.

\subsection{C3 — Active learning (planned)}
\label{subsec:c3}

Let the protocol $\Pi$ be \emph{chosen} by the learner: take infimum
over $\Pi$ in Theorem~\ref{thm:lower-bound}, \ref{thm:saturation}.
Theorem~\ref{thm:saturation} gives, for the optimal $\Pi^\star :=
\arg\max_\Pi \sigma_\Pi(\Scal_\Acal)$,
\begin{equation}
    L_{T^c}(\hat f_{\Acal^\star})
    \;\leq\; (1-\Rspec)\Var(\by) + O\left(\frac{\rho}{\sqrt{n \cdot \max_\Pi \sigma_\Pi(\Scal_\Acal)}}\right).
\end{equation}
Maximising $\sigma_\Pi(\Scal_\Acal)$ over sampling distributions is the
\emph{leverage-score sampling} problem on $\bB$, solved by sampling
$v \in V$ with probability proportional to $\|\bB(v, :)\|^2$
(Drineas--Mahoney 2018 and earlier).  This is the spectral-graph version
of A-optimal design; the resulting active-learning protocol is the one
that saturates Theorem~\ref{thm:lower-bound} fastest.  Full corollary
deferred to Paper 3.

% =========================================================================
\section{Open questions and remarks}
\label{sec:open}

\begin{enumerate}
\item \textbf{Sharpness of the saturation conjecture.}
Theorem~\ref{thm:saturation} bounds the Pesenson-ridge estimator above
by $(1-\Rspec)\Var + O(\rho/\sqrt{n\sigma_\Pi})$.  Is this slack the
\emph{minimax} excess loss, or is there a tighter estimator achieving,
e.g., $O(\rho/(n\sigma_\Pi))$?  The classical Cram\'er--Rao bound is
saturated with slack $0$ for efficient estimators in regular parametric
families; the graph-supervised analog of "regular parametric family" is
not yet pinned down.

\item \textbf{Non-ERM learners.}
Assumption~\ref{a:erm} (strict ERM) can be relaxed to \emph{approximate
ERM}: $L_T(\hat f_\Acal(T)) \leq \min_{f \in \Hcal_d} L_T(f) +
\epsilon_T$.  The bound \eqref{eq:thm-lower} then carries an additional
slack of $(\rho - 1) \cdot \E_T[\epsilon_T]$.  For regularised learners
(ridge, GBM with early stopping) this slack is typically $O(d/n)$ —
small but non-zero.

\item \textbf{Non-uniform protocols.}
The exchangeability hypothesis $\E_T[L_T(f)] = L_V(f)$ for deterministic
$f$ is exact for $\Pi_{\mathrm{LSO}}$ (uniform without replacement) and
$\Pi_{\mathrm{leave-1-out}}$, and approximate for GroupKFold.  For a
genuinely non-uniform protocol (e.g., active learning with leverage-score
sampling), this hypothesis is replaced by an importance-weighted version
$\E_T[\sum_{v\in T} w_T(v) \cdot (f(v)-y_v)^2] = L_V(f)$, which carries
through the proof with $L_T$ replaced by its weighted variant.

\item \textbf{Connection to graph Fisher information.}
Remark~\ref{rem:cramer-rao} interprets $\sigma_\Pi(\Scal_\Acal)$ as a
graph-Fisher-information analog: it controls the sharpness with which
the protocol identifies signals living in $\Scal_\Acal$.  A formal
definition of "graph Fisher information" — analogous to the classical
$I(\theta) = -\E[\partial^2_\theta \log p(\cdot;\theta)]$ — would
complete the analogy.  Tanaka--Eldar 2020 mention but do not develop
this connection.

\item \textbf{Saturation when $\bP_{\Scal_\Acal}\by \notin \Hcal_d$.}
Assumption~\ref{a:realisable} is the only place where we restrict
$\Hcal_d$.  When it fails (e.g., hard-threshold classifiers, decision
stumps of depth 0), the witness predictor in
Eq.~\eqref{eq:step3-ERM} is unavailable and the bound \eqref{eq:thm-lower}
weakens to a $\min_{f \in \Hcal_d} L_V(f)$-floor argument with no
guaranteed tightness.
\end{enumerate}

% =========================================================================
\section*{Acknowledgements (working notes)}

These notes formalise a research direction discussed in May 2026
between R.F.\ and an exploratory conversation with Claude
(Anthropic) on the trajectory of the cycling-data-lab structural-bounds
program.  The bibliographic verification of Remark~\ref{rem:pesenson-extension}
(R1: arbitrary-subspace Pesenson extension) confirmed that the
Tikhonov reconstruction bound on a generic feature subspace is already
in the published literature (Chepuri \& Leus 2018; Tanaka \& Eldar 2020),
so Lemma~\ref{lem:pesenson-subspace} above can be cited rather than
re-proven in any submission-ready manuscript.

\end{document}
